\section{The triangular lattice forms below \texorpdfstring{$\alpha_L = 0.372$}{alpha\_L = 0.372}}
\label{sec:5}

\subsection{Gravity adds an attractive channel to the vortex interaction}
\label{sec:5.1}

Below the critical coupling the normal phase is unstable, and we ask which
lattice forms and whether it is stable. We follow the probe-limit route
of~\cite{Bu:2012mq} and expand the free energy to quartic order in the
condensate, now on the brane and with the condensate sourcing the metric as
well as the abelian field. Just above the onset
$W_x = \eta\,w_0(r)\,\psi(x,y)$, $W_y = -iW_x$, with $w_0$ the zero mode of
section~\ref{sec:3} and $\psi$ a lowest-Landau-level function with one flux
quantum per unit cell. The Fourier coefficients $\omega_G$ of $|\psi|^2$,
with $G$ a reciprocal lattice vector, have $|\omega_G|^2 = e^{-\gamma_G/2}$ with
$\gamma_G = G^2/B$ (appendix~\ref{app:probe-lll-identity}), which depends
only on the shape of the lattice; we write $\gamma = s/B$ for a general
squared harmonic $s = G^2$. The free energy per unit area is
\begin{equation}
F = -\epsilon\,\eta^2 + \Big[\tfrac12 K_{G=0} + C(\tau;\beta)\Big]\eta^4~,
\qquad
C(\tau;\beta) = \tfrac12\sum_{G\neq0} e^{-\gamma_G/2}\,K(\gamma_G B_c(\beta);\beta)~,
\end{equation}
where $\epsilon \propto B - B_c$ and $\tau$ labels the lattice shape
(section~\ref{sec:5.2}). The uniform term $K_{G=0}$ is the same for every
lattice, and the shell sum $C$ selects the lattice.

The shell sum is built from the quartic kernel $K(s;\beta)$, the self-interaction
of the condensate through a density modulation of wavenumber $G$. The
density sources the abelian field and the metric, and since the background
depends on $r$ only, the response at each $G$ is a radial boundary value
problem depending on $s$ alone (appendix~\ref{app:response}). The kernel is
\begin{equation}
K(s;\beta) = \big(C_4 - X\big)(s;\beta) - \tfrac12\,\mathcal{P}(s;\beta)~, \qquad
C_4 = \int p_1\,w_0^4\,dr~,
\end{equation}
with $p_1 = e^{Z-2V}$ and $p_2 = Ue^Z$ the two weights of the
Sturm--Liouville problem of section~\ref{sec:3}, $(p_2 w')' + Bp_1 w = 0$.
We normalise the zero mode on each brane by $\int p_1 w_0^2\,dr = 1$. The
contact term $C_4$ is the self-repulsion of the condensate and $X$ its
screening by the induced $U(1)$ field. In the probe limit this bracket is
the whole kernel, and its positivity is what makes the vortices repel and
form a lattice. The new term $\mathcal{P}$ is the exchange of the metric
response between two insertions of the condensate stress tensor. Its
weight $\tfrac12$ relative to $X$ is fixed by the quadratic on-shell action
(appendix~\ref{app:response-metric}). We write $X_{\rm ab}$ for the
screening by the photon alone with the metric frozen; it equals $X$ in the
probe limit (appendix~\ref{app:response-abelian}). Pairing the metric
response with the photon's Maxwell stress as well gives the dressed
pairing $\mathcal{P}_{\rm dress}$, and $K = C_4 - X_{\rm ab} -
\tfrac12\mathcal{P}_{\rm dress}$ splits the kernel into gauge-invariant
parts (appendix~\ref{app:response-metric}). On the grid
$0 < \beta \le 45$ described below and on the throat,
$\mathcal{P}_{\rm dress} > 0$: gravity lowers
the kernel below the photon-only kernel $C_4 - X_{\rm ab}$, as the
screening lowers it below $C_4$.

The uniform term $K_{G=0}$ is the response to a strictly uniform
condensate at fixed temperature and field. It is a separate radial problem,
in which the photon does not respond, since the Bianchi identity fixes the
uniform part of $F_{xy}$ at $B$. We computed it directly and find that it
equals the long-wavelength limit of the kernel,
\begin{equation}
K_{G=0}(\beta) = K_0(\beta) \equiv \lim_{s\to0}K(s;\beta)~,
\end{equation}
to $10^{-8}C_4$ on every rung of the ladder up to $\beta = 10^8$; we write $K_0$ for both from here on. The agreement is not
automatic, because the gauge used at finite $G$ becomes singular as
$G \to 0$; appendix~\ref{app:response-metric} shows that the pairing is
insensitive to that singular gauge transformation. The quartic coefficient of a uniform lattice is
therefore $\tfrac12K_0 + C$.

We computed the kernel on a dense grid up to $\beta = 45$ ($\alpha = 0.1745$), at selected harmonics along the ladder up to its rung
$\beta = 10^8$ ($\alpha = 0.544$), and at $\alpha_\star$
as a fixed problem on the throat (appendix~\ref{app:numerics-scan}). On the
grid it is positive everywhere but suppressed at long wavelength: $K/C_4$
at $s = 0.02B_c$ falls from $0.98$ at $\beta = 0$ to $0.18$ at $\beta = 45$,
and from $\beta \approx 8$ the kernel has a shallow interior maximum at
$s \lesssim B_c$ (figure~\ref{fig:kernel}). At $\alpha_\star$ it is negative
at long wavelength.

\begin{figure}[t]
\centering
\includegraphics{kernel_family}
\caption{The quartic kernel. Left: $K(s;\beta)/C_4$ against $s/B_c$ for
$\beta = 0, 1, 4, 8, 16, 45$ (light to dark), with the first shells $2\pi$
(square) and $4\pi/\sqrt3$ (triangular) marked. Right: the throat kernel
$K/C_4$ at $\astar$ against $\gamma = s/B_c$,
negative below $\gamma_0 = 0.8465$ (shaded), with (inset) its gauge-invariant
split (appendix~\ref{app:response-metric}): the photon pairing
$(C_4 - X_{\rm ab})/C_4$ with the metric frozen, and the dressed metric
pairing $-\mathcal{P}_{\rm dress}/2C_4$, whose near-cancellation
produces it.}
\label{fig:kernel}
\end{figure}

\subsection{The triangular lattice is the best Bravais lattice up to \texorpdfstring{$\alpha_I$}{alpha\_I}}
\label{sec:5.2}

Up to rotation and scale, a one-flux-quantum Bravais lattice is fixed by a
modular parameter $\tau = \tau_1 + i\tau_2$ in the upper half plane, with
reciprocal harmonics
\begin{equation}
\gamma_{mn} = \frac{2\pi}{\tau_2}\,|m + n\tau|^2~, \qquad (m,n) \neq (0,0)~.
\end{equation}
The set of harmonics is invariant under $\tau \to \tau + 1$ and
$\tau \to -1/\tau$, so $C$ is modular invariant for any kernel. The
triangular lattice $\tau = e^{i\pi/3}$ and the square lattice $\tau = i$,
the elliptic fixed points, are therefore stationary at every $\beta$.

In the probe limit the kernel is completely monotone in $s$, and
Montgomery's
theorem~\cite{Montgomery:1988mt,Cohn:2007universally,Betermin:2016theta}
makes the triangular lattice the global minimum over all one-flux-quantum
Bravais lattices (appendix~\ref{app:probe}). On the brane we know of no
Stieltjes representation for the metric pairing, so the proof of
appendix~\ref{app:probe} does not carry over. The function the theorem
needs, $e^{-\gamma/2}K(\gamma B_c)$, is in fact not completely monotone from
$\beta \approx 4$, where it turns concave at long wavelength, and from
$\beta \approx 8$ the kernel itself is not monotone. This does not by itself disfavour the triangular
lattice~\cite{Betermin:2016theta,Betermin:2019noncm}, which fully
backreacted holographic lattices in a magnetic field also
prefer~\cite{Donos:2015eew,Donos:2020viz}.
We scanned $C$ over the fundamental domain up to $\tau_2 = 5$ on every
brane of the grid and every rung of the ladder up to $\beta = 10^8$
($\alpha = 0.544$; appendix~\ref{app:numerics-scan}). On this grid the minimiser never leaves the triangular
point, and the Hessian there stays positive and
isotropic. We measure the selection by the relative margin over the square
lattice, $(C_\square - C_\triangle)/C_\triangle$, which falls slowly from
$0.246$ in the probe limit to $0.227$ at $\alpha = 0.1745$ and $0.18$ at
$\beta = 10^8$. The shape term itself is small,
$C_\triangle \approx 0.01\,C_4$, so in the probe limit this margin is a
$0.56\%$ difference in condensation energy.

Towards the stripe edge $C$ grows in proportion to $\sqrt{\tau_2}$ times a
Gaussian average $I(\beta)$ of the kernel over momentum transfers of order
$B_c$ (appendix~\ref{app:numerics-scan}). This average stays
positive up to $\alpha_I = 0.494$, so for $\alpha < \alpha_I$, $C$ grows
towards the stripe edge, and with the edges scanned to $\tau_2 = 50$
(appendix~\ref{app:numerics-scan}) the triangular minimum is global. Above $\alpha_I$,
$C$ is unbounded below towards the stripe edge, a sign that the truncation
has left its domain of validity rather than a prediction of infinite
elongation. Since $\alpha_I > \alpha_L$, this only concerns lattices that
are already unstable (section~\ref{sec:5.3}).

At $\alpha_\star$ the selection becomes a $\beta$-independent problem on the
BTZ throat, because the frozen flux plane holds $s/B$ fixed
(appendix~\ref{app:response-throat}). The triangular lattice
still beats the square, with relative margin $+0.1766$, but the kernel is negative for
$\gamma < \gamma_0 = 0.8465$, because the lighter of the two neutral
channels of the metric response has a longer range than the photon that
screens the repulsion. The brane kernel on the ladder converges onto these
values (appendix~\ref{app:numerics-scan}).

\subsection{The uniform lattice destabilises at \texorpdfstring{$\alpha_L = 0.372$}{alpha\_L = 0.372}}
\label{sec:5.3}

Along the ladder $K_0/C_4$ falls monotonically and changes sign at
$\alpha = 0.3506$, $\beta = 3305$ (figure~\ref{fig:density}).
Beyond this coupling the interaction between density modulations of the
condensate is attractive at long wavelength. In the boundary theory this
is the exchange of the plasma's energy and pressure fluctuations between
the modulations. On the ladder $K_0/C_4$ saturates near $-0.06$, the value
of the throat kernel as $\gamma \to 0$ (figure~\ref{fig:kernel}, right).

A long-range attraction does not by itself destabilise the lattice, because
modulating the condensate density also modulates the lattice harmonics,
whose local energy resists it. Modulating the density of the uniform
lattice as $1 + \delta\cos(q\cdot x)$ on scales long compared with the
magnetic length adds harmonics at $G \pm q$ as well as at $q$, and to
second order the quartic energy rises by
\begin{equation}
\Delta F = \tfrac14\,\delta^2\eta^4 \sum_G e^{-\gamma_G/2}\,K\big(|G+q|^2;\beta\big)
\;\longrightarrow\; \tfrac14\,\delta^2\eta^4\,\big[K_0 + 2C_\triangle\big]
\qquad (q \to 0)~.
\end{equation}
The inverse compressibility of the lattice at fixed flux is thus
$K_0 + 2C_\triangle$, not $K_0$. To check that no other perturbation goes
first, we computed the second variation of the quartic functional about the
uniform lattice over all lowest-Landau-level perturbations at one flux
quantum. It splits over the magnetic Brillouin zone into $2\times2$ blocks
(appendix~\ref{app:response-hessian}). As $q \to 0$ the lower
branch is the Goldstone mode of translations and the upper branch tends to
$2(K_0 + 2C_\triangle)$. On a dense grid over the zone, from the probe
limit to $\alpha_L$, no eigenvalue is negative (in the probe limit the
lowest eigenvalue on the zone boundary is $0.0186\,C_4$, at the midpoint
$M$ of an edge). The
first instability of the uniform triangular lattice is therefore the
long-wavelength density mode, at $K_0 + 2C_\triangle = 0$, where the
quartic coefficient $\tfrac12K_0 + C_\triangle$ changes sign, that
is
\begin{equation}
\alpha_L = 0.3722
\end{equation}
(table~\ref{tab:couplings}).

Since $K_{G=0} = K_0$, the same combination is the quartic coefficient of
the uniform lattice, so $\alpha_L$ is also where the transition changes
order: it is a tricritical point of the onset line. It plays the role of
$\kappa = 1/\sqrt2$~\cite{Abrikosov:1956sx} in Ginzburg--Landau theory with a dynamical gauge
field, where the sign of the same Abrikosov quartic coefficient at
$H_{c2}$ separates type II from type I. Below $\alpha_L$ the normal phase gives way at $B_c$, in a
second-order transition, to a stable triangular lattice; between
$\alpha = 0.351$ and $\alpha_L$, where $K_0 < 0$, it is held stable by the
$2C_\triangle$ term. Above $\alpha_L$
the quartic coefficient is negative: the transition is first order, $B_c$
is the spinodal of the normal phase, and the uniform lattice is itself
unstable to long-wavelength modulation. The coefficient stays negative as
far as we computed the kernel, with
$(\tfrac12K_0 + C_\triangle)/C_4 = -0.025$ at $\beta = 10^8$, and what
replaces the lattice there is outside the quartic expansion. No other
one-flux-quantum Bravais lattice is stable above $\alpha_L$ either. Less
compact lattices keep their density stiffness $K_0 + 2C(\tau)$ longer; the
square keeps it to $\alpha = 0.3767$. But the square is a saddle of $C$ and
is unstable at the zone boundary at every coupling, and other shapes are
not stationary. Above $\alpha_L$, every lattice that still has positive
density stiffness is unstable to shear (appendix~\ref{app:numerics-scan}).\footnote{The electric
$p$-wave transition of~\cite{Ammon:2009xh} also turns first order beyond a
coupling, $\alpha \approx 0.13$ in our variable.}

\begin{table}[t]
\centering
\begin{tabular}{lccc>{\raggedright\arraybackslash}p{0.40\textwidth}}
\toprule
coupling & $\beta$ & $B_c u_H^2$ & $T_c/T_c^{\rm probe}$ & what changes there \\
\midrule
$0.351$ & $3.31\times10^3$ & $97.1$ & $0.23$ & long-wavelength vortex interaction turns attractive, $K_0 = 0$ \\
$\alpha_L = 0.372$ & $6.51\times10^3$ & $132.3$ & $0.20$ & uniform lattice unstable; transition first order, $B_c$ a spinodal \\
$\alpha_\star = 2/3$ & $\infty$ & $\infty$ & $0$ & no linear instability of the charged sector at any $B$, $T$ \\
\bottomrule
\end{tabular}
\caption{The couplings that organise the lattice and the onset, with the
backreaction parameter and the critical field at fixed temperature on the
onset curve, and the onset temperature at fixed field relative to its probe
value, $T_c/T_c^{\rm probe} = (5.13/B_c u_H^2)^{1/2}$. Below $\alpha_L$ a
triangular vortex crystal forms in a second-order transition; whether a
condensed phase exists above $\alpha_\star$ is open.}
\label{tab:couplings}
\end{table}

\begin{figure}[t]
\centering
\includegraphics{density_threshold}
\caption{The long-wavelength kernel along the ladder: $K_0/C_4$ (circles)
and the quartic coefficient of the uniform lattice
$(\tfrac12 K_0 + C_\triangle)/C_4$ (squares) on the rungs $\beta = 1$ to
$10^8$ against $\alpha$, with $\ln\beta$ of the ladder on the top axis. The
long-range interaction turns attractive where $K_0$ changes sign, at
$\alpha = 0.351$; the uniform lattice destabilises, and the transition
turns first order, where the quartic coefficient does, at
$\alpha_L = 0.372$. The crystal and first-order regimes, and the regime
without a charged instability, are shaded.}
\label{fig:density}
\end{figure}

At fixed $\alpha$ the only scale is the temperature, so the onset is
$B_c = f(\alpha)\,T^2$ exactly, with $f = \pi^2 B_c u_H^2$ from
figure~\ref{fig:1}. Cooling at fixed field therefore meets it at
$T_c = \sqrt{B/f(\alpha)}$ for every $\alpha < \alpha_\star$. Gravity lowers
$T_c$ to a fifth of its probe value by $\alpha_L$
(table~\ref{tab:couplings}) and to under two percent of it at
$\alpha = 0.55$, and $T_c$ vanishes with the essential singularity of
section~\ref{sec:4.6} at $\alpha_\star$.
